% palette.tex — Agentra house palette + style language for draw-figure-dph.
% SINGLE SOURCE OF TRUTH. Figures must \input this, not re-declare colors.
%
% Identity: near-monochrome forest ink, agentra-green as the ONLY structural
% color, ONE warm-orange highlight. NOT a 6-color rainbow.
%
% ROLE IS ENCODED TWICE: by hue AND by rule line-style. Hue alone fails in
% grayscale (agreen luma 134, warm 142, g500 113 all collapse), so every role
% rule also carries a distinct dash pattern. Never encode a role by hue alone.
%
% Required libs: positioning, shapes.geometric, arrows.meta, fit, backgrounds, calc

% ---- brand colors ----
\definecolor{forest}{HTML}{0C2F1F}  % ink: ALL text + strokes
\definecolor{agreen}{HTML}{46B152}  % structural: stages, spine, badges
\definecolor{lime}{HTML}{B3D335}    % secondary track (e.g. attacker/auxiliary pipeline)
\definecolor{warm}{HTML}{D97742}    % the SINGLE warm highlight: decision + output/alert ONLY
\definecolor{g500}{HTML}{6B7280}    % neutral parameter / aux
\definecolor{g300}{HTML}{DADEE2}    % hairlines, group borders

% ---- knobs. Defaults ARE the canonical house values (examples/ex3b-agentra-style).
% Override by \def-ing BEFORE \input; only do so for slides.
\providecommand{\ruleW}{1.6pt}      % role top-rule thickness   (slides: 2.4pt)
\providecommand{\cardStroke}{0.5pt} % card border               (slides: 0.6pt)
\providecommand{\cardW}{32mm}       % uniform card width; the diamond tracks it
% The decision diamond's stroke and proportion live in knobs so the `dec` node and its
% legend swatch \swDec cannot drift apart. \swDec used to hard-code aspect=1.6 and
% 0.7pt (plus a minimum height, which overrides the aspect entirely), so the key showed
% a visibly different diamond from the diamonds it was explaining.
\providecommand{\decLW}{1.1pt}      % decision diamond stroke   (slides: 1.6pt)
\providecommand{\decAspect}{2.2}    % decision diamond width:height
\providecommand{\railW}{1.0pt}      % legacy (spine motif retired), kept for old files
\providecommand{\badgeSz}{4.4mm}    % legacy (badge motif retired)

% ---- style language: white cards, hairline forest border, role via TOP-RULE ----
\tikzset{
  card/.style={draw=forest,line width=\cardStroke,fill=white,text=forest,align=center,
    rounded corners=1pt,inner sep=5pt,minimum height=10mm,minimum width=\cardW},
  amod/.style={draw=forest,line width=\cardStroke,fill=white,text=forest,align=center,
    rounded corners=1pt,inner sep=5pt,minimum height=10mm,minimum width=34mm},
  % diamond width tracks \cardW so it stays in the same column grid as the cards
  dec/.style ={diamond,aspect=\decAspect,draw=warm,line width=\decLW,fill=white,text=forest,
    align=center,inner sep=1pt,minimum width=\cardW},
  % panel: labeled region grouping one actor's blocks. Panels are what give the
  % figure structure and density -- prefer them over a long vertical rail.
  panel/.style ={draw=g300,line width=0.9pt,rounded corners=2pt,inner sep=6pt},
  ptitle/.style={font=\scriptsize\bfseries,text=forest,inner sep=1.5pt,fill=white},
  % step chip: small number sitting ON the role rule (replaces the old spine+badge motif)
  % \scriptsize, never \tiny: at 10pt base \tiny is 5pt, which lands at 4.6-5.0pt once
  % the artwork is placed. The step number is real information, not decoration.
  chip/.style={circle,fill=agreen,text=white,font=\scriptsize\bfseries,
    minimum size=4.4mm,inner sep=0pt},
  badge/.style={circle,fill=agreen,draw=forest,line width=0.5pt,text=white,
    font=\footnotesize\bfseries,minimum size=\badgeSz,inner sep=0pt},
  data/.style  ={->,draw=forest,line width=0.8pt,shorten >=1mm,shorten <=1mm},
  active/.style={->,draw=lime!72!forest,line width=1.2pt,dotted,shorten >=1mm},
  param/.style ={->,draw=g500,line width=0.8pt,dashed,shorten >=1mm,shorten <=1mm},
}

% ---- ROLE RULES: hue + dash pattern (both channels). Use these, never a raw \draw. ----
\newcommand{\ruleStage}[1]{% always-executed pipeline stage
  \draw[agreen,line width=\ruleW] (#1.north west)--(#1.north east);}
\newcommand{\ruleTrack}[1]{% secondary track (attacker / auxiliary pipeline)
  \draw[lime!80!forest,line width=\ruleW,dash pattern=on 3.5pt off 2.2pt]
    (#1.north west)--(#1.north east);}
\newcommand{\ruleAux}[1]{% neutral parameter / constraint block
  \draw[g500,line width=\ruleW,dash pattern=on 1pt off 2pt]
    (#1.north west)--(#1.north east);}
% A `double` stroke is 2*line width + double distance wide. Keep that TOTAL equal to
% \ruleW, or the output rule renders ~2.6x heavier than the stage rules and its white
% core erases the card's own top border.
\providecommand{\ruleOutLW}{0.5pt}\providecommand{\ruleOutDD}{0.6pt} % DECK: 0.9pt / 0.8pt
\newcommand{\ruleOut}[1]{% decision result / output / alert  (split bar = unmistakable in print)
  \draw[warm,line width=\ruleOutLW,double,double distance=\ruleOutDD]
    (#1.north west)--(#1.north east);}
% ---- LEGEND SWATCHES: always build the key from these, never hand-draw it. ----
% A hand-drawn key hard-codes widths/dash patterns and silently desynchronises from
% the real rules the moment \ruleW (or the slide override) changes.
\newcommand{\swStage}{\tikz\draw[agreen,line width=\ruleW](0,0)--(0.38,0);}
\newcommand{\swTrack}{\tikz\draw[lime!80!forest,line width=\ruleW,
  dash pattern=on 3.5pt off 2.2pt](0,0)--(0.38,0);}
\newcommand{\swAux}{\tikz\draw[g500,line width=\ruleW,
  dash pattern=on 1pt off 2pt](0,0)--(0.38,0);}
\newcommand{\swOut}{\tikz\draw[warm,line width=\ruleOutLW,double,
  double distance=\ruleOutDD](0,0)--(0.38,0);}
\newcommand{\swData}{\tikz\draw[data,shorten >=0pt,shorten <=0pt](0,0)--(0.38,0);}
\newcommand{\swActive}{\tikz\draw[active,shorten >=0pt](0,0)--(0.38,0);}
\newcommand{\swParam}{\tikz\draw[param,shorten >=0pt,shorten <=0pt](0,0)--(0.38,0);}
% A diamond with EMPTY text has zero natural height, so `aspect` alone leaves the
% swatch a flat sliver: the height must be given. It is therefore COMPUTED from
% \decAspect rather than hard-coded, so the swatch keeps the real node's proportion
% (the old 5.5mm x 3mm at aspect=1.6 drew a 1.83 ratio against the node's 2.2).
\providecommand{\swDecW}{5.5mm}
\newcommand{\swDec}{\tikz{%
  \pgfmathsetlengthmacro{\swDecH}{\swDecW/\decAspect}%
  \node[draw=warm,line width=\decLW,fill=white,diamond,aspect=\decAspect,
        minimum width=\swDecW,minimum height=\swDecH,inner sep=0pt]{};}}

% step number sitting on the rule: \stepchip{node}{n}
\newcommand{\stepchip}[2]{\node[chip] at ([xshift=5mm]#1.north west){#2};}
% panel title: \paneltitle{panel node}{text}
\newcommand{\paneltitle}[2]{\node[ptitle,anchor=west] at ([xshift=4mm]#1.north west){#2};}
%
% NOTE: do NOT name these \toprule — that collides with booktabs in a real paper preamble.
% NOTE: never name a TikZ style `out`, `in`, `step`, `scale` — they shadow keys.
% DECK intensity, before \input:
%   \def\ruleW{2.6pt}\def\railW{1.8pt}\def\badgeSz{6.2mm}\def\cardStroke{0.6pt}\def\cardW{48mm}
